% Figure 3 — where the workspace comes from, and what is kept of it.
%
% The agent is not handed a clone: it is handed an extracted archive
% committed once, so the tree carries no history to read the answer out
% of. What comes back out is that tree plus one diff — read by host-side
% git from a `.git/` the agent could write to.
%
% Two rows, joined at the end of the first: preparing the tree is one
% chain, capturing the work is another, and they meet where the agent
% does.
%
% Sources: src/bench/solves.py materialise / capture /
% _refuse_stack_files. The two probes that made this necessary: a clone
% carrying 3,671 commits, and a worktree that shared every ref.
\documentclass[tikz,border=4pt]{standalone}
\input{style}

\begin{document}
\begin{tikzpicture}

% --------------------------------------------------------------- row one --
\node[commit, text width=2.4cm] (pin) at (1.9,10)
  {\icgit\ the pin\\ \texttt{089d2531}};
\node[code, text width=2.7cm, align=center] (arch) at (5.2,10)
  {git archive\\<pin>};
\node[code, text width=3.1cm, align=center] (extract) at (8.8,10)
  {tar -x -C\\out/workspaces/\\\textit{run}/};
\node[warn, text width=2.6cm, align=center] (refuse) at (12.2,10)
  {\icno\ \textbf{refuse}\\ stack files};
\node[code, text width=3.1cm, align=center] (baseline) at (15.8,10)
  {git init\\add -A\\commit "baseline"};

\draw[flow] (pin) -- (arch);
\draw[flow] (arch) -- (extract);
\draw[flow] (extract) -- (refuse);
\draw[flow] (refuse) -- (baseline);

\node[note, text width=2.8cm, align=center] at (5.2,8.75)
  {no clone:\\no history, no refs};
\node[note, text width=3.6cm, anchor=north west] at (16.0,9.0)
  {one commit · no remote\\ no refs · no reflog\\
   author \texttt{bench <bench@localhost>}};

% --------------------------------------------------------------- row two --
\node[device, text width=3.2cm] (agent) at (6.2,5.4)
  {\icagent\ \textbf{the agent}\\ edits in place, as root,\\
   through the bind mount};
\node[code, text width=3.3cm, align=center] (diff) at (10.6,5.4)
  {git add -N .\\git diff --binary\\\normalfont\footnotesize on the host};
\node[file, text width=3.4cm] (out) at (14.9,5.4)
  {\icvolume\ \textbf{out/solves/\textit{run}/}\\ \texttt{record.json}\\
   \texttt{diff.patch} \quad \texttt{tree/}};

\draw[flow] (agent) -- (diff);
\draw[flow] (diff) -- (out);

\node[note, text width=3.3cm, align=center] at (11.1,3.95)
  {intent-to-add is what\\ puts new files in the patch};
\node[note, text width=3.4cm, align=center] at (14.9,3.95)
  {the tree is copied aside\\ as it was left};

% --------------------------------------------------------------- the join --
\draw[live] (baseline.south) |-
  node[lbl, above, pos=0.7] {the tree} (agent.north);

% the .git the agent can write, and the git that reads it
\node[breach, text width=3.2cm, anchor=north] (dotgit) at (6.2,3.9)
  {\icgit\ \texttt{.git/} is inside the\\mount: config, hooks,\\
   objects are the agent's};
\draw[attack] (dotgit.north east) -- (diff.south west);

% ---------------------------------------------------------------- caveat --
\node[warn, anchor=north west, text width=15.6cm] at (1.0,2.0)
  {\icwarn\ \textbf{What the patch cannot show.} The diff is read through
   git, so a write to a gitignored path is invisible: \texttt{diff\_bytes:
   0} means \emph{not measured}, never \emph{changed nothing}. And the git
   that reads it runs on the host, as the supervisor, against a repository
   the agent could rewrite: a \texttt{core.fsmonitor} line in
   \texttt{.git/config} is a command (figure 5).};

\end{tikzpicture}
\end{document}
